% Overleaf: upload this file and figures/a/throughput.pdf and
% figures/a/a0_latency_share.pdf, preserving the figures/a/ folder structure.
% Prefer PDF: these charts are vector graphics and stay sharp when scaled.
% PNG is also supported, but is a raster fallback.
% This section uses the existing report macros \code, \expt, \experiment,
% and \resultfigure, and requires the booktabs and graphicx packages.

\section{Task 1: Baseline Design, Workload A}\label{sec:task1}

\subsection{Experimental Setup}\label{sec:setup}

The experiments run \code{run.sh} under Slurm on the SoC compute cluster,
using Citus 14.1.0 on PostgreSQL 18.6 with 4 CPUs allocated per node.
One node acts as coordinator and runs \code{pgbench}; four nodes act as
workers, except in \expt{a0\_1}, which uses two workers.
The recorded \code{topology.txt} files identify \code{xcnf1} as coordinator
and \code{xcnf3}--\code{xcnf6} as workers for \expt{a0};
\expt{a0\_1} uses coordinator \code{xcnf9} and workers
\code{xcnf10}--\code{xcnf11}; and \expt{a0\_2} and \expt{a0\_3}
use coordinator \code{xcnf9} and workers \code{xcnf10}--\code{xcnf13}.
Consequently, every comparison with \expt{a0} also changes physical nodes.
Furthermore, the three parameter experiments have overlapping execution
times on shared nodes. Their results describe the observed sensitivity
of the baseline, but do not isolate the causal effect of each parameter.

\begin{table}[htbp]
  \centering
  \begin{tabular}{@{}p{0.30\linewidth}p{0.65\linewidth}@{}}
    \toprule
    Parameter & Value \\
    \midrule
    Client counts & 1, 4, 8, 16, 32, 64 \\
    Configured measured run / warm-up & 600\,s / 60\,s per client count \\
    Isolation level & \code{SERIALIZABLE} \\
    Retries & \code{pgbench -{}-max-tries=3} \\
    Baseline dataset (\code{SCALE=1}) &
      1,000 merchants, 5,000,000 orders, approximately 15,000,000 order items \\
    Baseline skew & 10\% of merchants receive about 50\% of order activity \\
    \bottomrule
  \end{tabular}
  \caption{Benchmark configuration. Worker count, skew, and scale vary in
  \expt{a0\_1}, \expt{a0\_2}, and \expt{a0\_3}, respectively.}
  \label{tab:setup}
\end{table}

Average throughput is the main comparison metric, accompanied by the
p50, p95, and p99 latencies of successful transactions.
All tables reproduce \code{summary.csv}, whose TPS is the number of
successful logged transactions divided by the configured 600\,s.
This convention includes transactions that finish while the run drains
after the configured duration, and therefore exceeds elapsed-time TPS
when draining is substantial. For example, at 64 clients the reported
TPS is 3.203 for \expt{a0} and 2.083 for \expt{a0\_3}, whereas
\code{pgbench} reports elapsed-time TPS of 2.951 and 1.768, respectively.
A client count is invalid whenever \code{pgbench} exits non-zero;
such results are retained as diagnostics but excluded from performance
conclusions. Invalid counts are 16 and 32 in \expt{a0}, and 32 in
\expt{a0\_2}, all with exit code 2.
A successful process exit does not imply that every transaction succeeded:
serialization failures after retries are reported separately.
Each configuration has one recorded run per client count, so the
comparisons do not establish statistical significance.

\experiment{a0}{Baseline Design}

The schema comprises merchants, customers, products, inventory, orders,
and order items: merchants remain local to the coordinator, while the
other five tables are distributed by their respective primary IDs
(\code{customer\_id}, \code{product\_id}, \code{inventory\_id},
\code{order\_id}, and \code{order\_item\_id}).
In addition to primary-key indexes, the design creates five single-column
\code{merchant\_id} indexes, one on each distributed table, without
aligning distribution keys with the merchant predicates or join keys.

\resultfigure{figures/a/throughput.pdf}{Workload A throughput for the baseline
and its parameter experiments. Hollow markers denote invalid runs:
\expt{a0} at 16 and 32 clients, and \expt{a0\_2} at 32 clients.
Lines connect all measurements, including invalid runs. TPS uses the configured
600\,s denominator; hardware and execution overlap limit causal comparisons.}{fig:a0-throughput}

\begin{table}[htbp]
  \centering
  \begin{tabular}{@{}rrrrrl@{}}
    \toprule
    Clients & TPS & p50 (ms) & p95 (ms) & p99 (ms) & Validity \\
    \midrule
     1 & 3.225 & 11.232 &   2639.466 &   2687.850 & Valid \\
     4 & 3.573 & 26.328 &   9837.977 &  10314.475 & Valid \\
     8 & 3.532 & 40.429 &  19512.201 &  21173.629 & Valid \\
    16 & 3.242 & 52.678 &  33081.735 &  36601.314 & Invalid \\
    32 & 3.640 & 51.937 &  65001.448 &  68832.439 & Invalid \\
    64 & 3.203 & 71.790 & 128909.918 & 135395.836 & Valid \\
    \bottomrule
  \end{tabular}
  \caption{Baseline results from \expt{a0}. Invalid rows contain incomplete
  results and are excluded from comparisons.}
  \label{tab:a0-results}
\end{table}

The baseline gains little throughput from additional clients.
The highest valid measurement is 3.573 TPS at four clients, only
10.8\% above the single-client result; at 64 clients, throughput is
3.203 TPS. Meanwhile, p99 increases from 2.688\,s to 135.396\,s,
although p50 remains below 72\,ms. This separation indicates that
the aggregate median conceals a small set of very expensive transactions.

To quantify their contribution, the share of summed successful-transaction
latency attributable to transaction type \(i\) is
\[
  s_i =
  \frac{n_i\overline{\ell}_i}
       {\sum_j n_j\overline{\ell}_j},
\]
where \(n_i\) is its successful transaction count and
\(\overline{\ell}_i\) its average latency.
This measures accumulated client latency, including waiting and
communication, rather than CPU time.
At one client, T3 (Customer Order History) accounts for 76.6\% of
summed latency despite representing only 9.0\% of successful transactions;
T4 (Product Lookup) contributes another 16.9\%.
At 64 clients, their shares are 66.7\% and 30.8\%, respectively,
together accounting for 97.5\%.
T3's average latency rises from 2.637\,s to 124.031\,s, and T4's
from 0.479\,s to 55.135\,s.

\resultfigure{figures/a/a0_latency_share.pdf}{Shares of summed
successful-transaction latency in \expt{a0} at one and 64 clients.
T3 and T4 dominate both measurements.}{fig:a0-latency-share}

The transaction SQL explains why these reads are costly:
T3 joins orders to order items on \code{order\_id}, although order items
are distributed by \code{order\_item\_id}; T4 joins products to inventory
on \code{product\_id}, although inventory is distributed by
\code{inventory\_id}.
The joins therefore cannot rely on aligned distribution keys.
A separate baseline \code{EXPLAIN ANALYZE} run corroborates this
interpretation: T3 uses two map/merge stages, each with 32 map tasks and
16 merge tasks, while T4 includes a 32-map-task stage despite selecting
a single product. Their measured execution times in that diagnostic run
are approximately 2.748\,s and 0.446\,s.
The latency shares and plans support distributed join processing as
the principal baseline limitation; increasing concurrency mainly
increases latency without providing sustained throughput growth.
Serialization failures are an additional constraint at higher concurrency:
the valid 64-client run records 62 such failures, while the aborted
16- and 32-client runs cannot establish performance at those counts.

\experiment{a0\_1}{Baseline with Two Workers}

This experiment reduces the worker count from four to two, keeping
\code{SCALE=1}, the 50\% hot-order share, the baseline schema, and the
transaction mix unchanged. The recorded topology is
\code{xcnf9}--\code{xcnf11}, so comparison with \expt{a0} also changes
hardware and is affected by overlapping jobs on shared nodes.

\begin{table}[htbp]
  \centering
  \begin{tabular}{@{}rrrrr@{}}
    \toprule
    Clients & TPS & p50 (ms) & p95 (ms) & p99 (ms) \\
    \midrule
     1 & 2.190 &   9.855 &   4355.126 &   4442.821 \\
     4 & 2.158 &  29.394 &  17024.607 &  17907.785 \\
     8 & 2.105 &  51.005 &  33676.193 &  37241.774 \\
    16 & 2.098 &  90.778 &  64270.391 &  71418.325 \\
    32 & 2.168 & 104.133 & 115890.853 & 126777.144 \\
    64 & 2.127 & 110.731 & 217663.624 & 229562.480 \\
    \bottomrule
  \end{tabular}
  \caption{Results for \expt{a0\_1}; all client counts have exit code 0.}
  \label{tab:a0-1-results}
\end{table}

Throughput remains between 2.098 and 2.190 TPS across client counts.
At the jointly valid counts of 1, 4, 8, and 64, it is
32.1--40.4\% below the baseline, while p99 is approximately
1.65--1.76 times higher.
T3 accounts for 82.0\% of summed latency at one client and 69.9\%
at 64 clients, with average latencies of 4.346\,s and 208.974\,s.
The results are consistent with reduced worker capacity making the
expensive distributed reads slower, but hardware changes and shared-node
contention prevent attributing the full difference to worker count.
There are 45 serialization failures at 64 clients despite a successful
process exit.

\experiment{a0\_2}{Baseline with Stronger Merchant Skew}

This experiment increases \code{HOT\_ORDER\_PERCENT} from 50 to 80.
The hot group remains 10\% of merchants, and the change applies both
to generated historical orders and to transactions that select a
merchant. Scale and worker count remain at one and four, respectively;
the run uses coordinator \code{xcnf9} and workers
\code{xcnf10}--\code{xcnf13}.

\begin{table}[htbp]
  \centering
  \begin{tabular}{@{}rrrrrl@{}}
    \toprule
    Clients & TPS & p50 (ms) & p95 (ms) & p99 (ms) & Validity \\
    \midrule
     1 & 3.127 & 11.398 &   2722.832 &   2783.007 & Valid \\
     4 & 3.360 & 34.812 &  10190.109 &  10669.536 & Valid \\
     8 & 3.397 & 58.757 &  19859.063 &  21515.363 & Valid \\
    16 & 3.203 & 74.209 &  37148.299 &  40242.646 & Valid \\
    32 & 3.497 & 66.868 &  64771.685 &  69502.688 & Invalid \\
    64 & 3.078 & 97.803 & 133177.781 & 140907.333 & Valid \\
    \bottomrule
  \end{tabular}
  \caption{Results for \expt{a0\_2}. The 32-client run exits with code 2
  and is excluded from performance comparisons.}
  \label{tab:a0-2-results}
\end{table}

At the jointly valid counts of 1, 4, 8, and 64, throughput is
3.0--6.0\% below \expt{a0}, and p99 is 1.6--4.1\% higher.
Median latency is more sensitive at concurrency: at eight clients it
increases from 40.429\,ms to 58.757\,ms, and at 64 clients from
71.790\,ms to 97.803\,ms.
Nevertheless, T3 and T4 still account for 97.4\% of summed latency
at 64 clients, indicating that the distributed read bottleneck persists.
Because the baseline distributes by entity IDs, stronger merchant skew
does not directly concentrate all hot-merchant data on one worker.
The observations do not establish that skew caused the differences:
physical nodes change relative to \expt{a0}, execution overlaps other
jobs, and no repeated measurements quantify run-to-run variation.
The valid 64-client run records 59 serialization failures.

\experiment{a0\_3}{Baseline with Double Data Scale}

This experiment increases \code{SCALE} from one to two while retaining
four workers and the baseline skew.
The initial dataset contains 2,000 merchants, 2,000,000 customers,
1,000,000 products and inventory rows, 10,000,000 orders, and
approximately 30,000,000 order items.
It uses coordinator \code{xcnf9} and workers
\code{xcnf10}--\code{xcnf13}, changing hardware relative to \expt{a0}
and overlapping the other parameter experiments during part of its run.

\begin{table}[htbp]
  \centering
  \begin{tabular}{@{}rrrrr@{}}
    \toprule
    Clients & TPS & p50 (ms) & p95 (ms) & p99 (ms) \\
    \midrule
     1 & 2.018 &  13.388 &   4657.780 &   4862.633 \\
     4 & 2.115 &  37.916 &  17199.926 &  18143.604 \\
     8 & 2.090 &  61.901 &  34566.551 &  36330.284 \\
    16 & 2.103 & 113.060 &  64815.498 &  71173.488 \\
    32 & 2.147 & 123.960 & 115263.450 & 128941.635 \\
    64 & 2.083 & 147.993 & 218813.640 & 228174.357 \\
    \bottomrule
  \end{tabular}
  \caption{Results for \expt{a0\_3}; all client counts have exit code 0.}
  \label{tab:a0-3-results}
\end{table}

Throughput remains between 2.018 and 2.147 TPS.
At the jointly valid counts of 1, 4, 8, and 64, it is
35.0--40.8\% below \expt{a0}, while p99 is approximately
1.69--1.81 times higher.
At one client, T3's average latency increases from 2.637\,s in the
baseline to 4.698\,s, and T6 (Global Sales) increases from
0.262\,s to 0.534\,s.
These increases are consistent with greater data-processing work,
particularly for distributed joins and aggregation.
At 64 clients, T3 and T4 account for 97.1\% of summed latency,
and overall p99 reaches 228.174\,s.
The run records 50 serialization failures at that client count.
The measurements show that the baseline remains dominated by expensive
reads at the larger scale, but the node change and overlapping execution
prevent interpreting the observed throughput reduction as an isolated
scale effect.
